% TOP_PROBLEM: {"id":30007752,"problem_number":"LOCAL-30007752","title":"Optimal income taxation with heterogeneous people","category_id":21,"category":{"id":21,"name":"economics","display_name":"Economics","description":"Open economic theory statements, published research agendas, and proposed scoped specifications; question provenance and historical priority are qualified per record.","slug":"economics","order_index":21,"created_at":"2026-10-08T00:00:00Z"},"status":"research_question","external_url":null,"legacy_ids":[],"published":false,"statement":"In the explicitly specified setting: How should income taxes account for heterogeneous preferences, earning abilities, bargaining power, and endogenous wages? The mathematical statement above fixes the target, assumptions and scope of this catalogue formulation.","economics":{"original_id":241,"field":"Taxation and social insurance","kind":"design","formalization_basis":"adapted_research_question","source_status":"research_agenda","priority_status":"earliest_found","priority_note":"Earliest explicit agenda verified in this audit, not a claim of first historical publication. The mathematical scope is authored for this catalogue.","earliest_verified_reference":{"authors":["Louis Kaplow"],"title":"Optimal Income Taxation","year":2022,"url":"https://laweconcenter.law.harvard.edu/wp-content/uploads/2024/11/Kaplow_1083.pdf","doi":"10.3386/w30199","locator":"Olin Discussion Paper 1083, July 2022; sections 3.2, 7, 8.2 and 9, printed pp. 89–90","evidence_summary":"The conclusion identifies further work on heterogeneous preferences, family composition, endogenous wages, and integrating tax instruments. This is an explicit research agenda, not the exact optimization model below."},"current_reference":{"authors":["Louis Kaplow"],"title":"Optimal Income Taxation","year":2024,"url":"https://www.aeaweb.org/articles?id=10.1257/jel.20221647","doi":"10.1257/jel.20221647","locator":"Journal of Economic Literature 62(2), 637–738; publisher abstract","evidence_summary":"The published survey describes heterogeneous preferences, families, endogenous wages and interactions with other instruments."},"formalization_note":"Catalogue-authored scoped formalization. Population, instruments, welfare objective and identification restrictions must be instantiated before claiming a resolution; source authors are not credited with these added assumptions."},"statusQualification":"Published question or research agenda verified at the cited locator. The mathematical specification is adapted for this catalogue and includes additional assumptions; source publication does not establish first-ever priority or certify this exact model remains unsolved. Catalogue-authored scoped formalization. Population, instruments, welfare objective and identification restrictions must be instantiated before claiming a resolution; source authors are not credited with these added assumptions.","statusReviewedAt":"2026-10-08"}
% FIRST_POSTED: 2026-10-08
% ENTRY_KIND: statement_only
% !TeX program = lualatex
\documentclass[11pt]{article}
\usepackage{fontspec}
\usepackage[a4paper,margin=25mm]{geometry}
\usepackage{amsmath,amssymb,mathtools}
\usepackage{xurl}
\usepackage[hidelinks]{hyperref}
\newcommand{\sref}[2]{\hyperlink{source:#1:#2}{[S#2]}}
\setlength{\emergencystretch}{3em}
\begin{document}
% problemId:30007752
\section*{Optimal income taxation with heterogeneous people}\label{30007752}
\subsection{Definitions and mathematical statement}
Fix a compact type space \(\Theta\), probability distribution \(F\), and occupations \(j=1,\ldots,J\). A household type \(\theta\) comprises abilities \(a_{\theta j}\), preferences and a specified intrahousehold bargaining weight. Let \(u_\theta(c,\ell,j)\) be its resulting welfare index, increasing in consumption \(c\geq0\) and decreasing in labor \(\ell\geq0\). A continuously differentiable production function determines competitive wages \(w_j\) from aggregate effective labor. For an income-tax schedule \(T\), household earnings are \(y=w_j a_{\theta j}\ell\), and households maximize \(u_\theta(y-T(y),\ell,j)\) over choices with \(y-T(y)\geq0\). Let \(\mathcal E(T)\) denote all allocations and wages satisfying household optimization, firm optimization and market clearing. Specify a welfare aggregator \(G_\theta\), revenue requirement \(R\), and admissible tax set \(\mathcal T\), including slope bounds and a treatment of zero income. Characterize
\[\sup_{T\in\mathcal T}\ \inf_{e\in\mathcal E(T)}\int_\Theta G_\theta(u_\theta(e))\,dF(\theta)\quad\text{subject to}\quad \int_\Theta T(y_\theta(e))\,dF(\theta)\geq R\ \text{for every }e\in\mathcal E(T).\]
The revenue constraint must hold in every equilibrium considered. Taxes with no equilibrium are inadmissible. Determine which joint distributions of abilities, preferences and bargaining weights are identified from available income, labor and family data, and how the optimal schedule varies across observationally compatible distributions. The target is robust policy characterization and empirically justified restrictions, rather than a universal numerical tax rate.

\par\textbf{Formulation and scope.} Catalogue-authored scoped formalization. Population, instruments, welfare objective and identification restrictions must be instantiated before claiming a resolution; source authors are not credited with these added assumptions.

\par\textbf{Open-question provenance.} An explicit question or research agenda was verified at the source locator. This does not by itself certify that every later version remains unresolved \sref{30007752}{1}.

\par\textbf{Historical priority.} Earliest explicit agenda verified in this audit, not a claim of first historical publication. The mathematical scope is authored for this catalogue.

\subsection{Short English statement}
In the explicitly specified setting: How should income taxes account for heterogeneous preferences, earning abilities, bargaining power, and endogenous wages? The mathematical statement above fixes the target, assumptions and scope of this catalogue formulation.

\subsection{Sources}
\begin{itemize}\raggedright
\item[\textnormal{[S1]}]\hypertarget{source:30007752:1}{} Louis Kaplow. 2022. Optimal Income Taxation. Olin Discussion Paper 1083, July 2022; sections 3.2, 7, 8.2 and 9, printed pp. 89–90.
\url{https://laweconcenter.law.harvard.edu/wp-content/uploads/2024/11/Kaplow_1083.pdf}.
DOI: 10.3386/w30199.
{\small\textit{Earliest verified question/agenda reference; historical priority qualified below. The conclusion identifies further work on heterogeneous preferences, family composition, endogenous wages, and integrating tax instruments. This is an explicit research agenda, not the exact optimization model below.}}

\item[\textnormal{[S2]}]\hypertarget{source:30007752:2}{} Louis Kaplow. 2024. Optimal Income Taxation. Journal of Economic Literature 62(2), 637–738; publisher abstract.
\url{https://www.aeaweb.org/articles?id=10.1257/jel.20221647}.
DOI: 10.1257/jel.20221647.
{\small\textit{Supporting or current literature; see evidence qualification. The published survey describes heterogeneous preferences, families, endogenous wages and interactions with other instruments.}}
\end{itemize}
\end{document}
